\let\stop\empty
\documentclass{exam-zh}
\usepackage{siunitx}
\usepackage{tkz-euclide}
\usepackage{diagbox}  
\usepackage{lastpage}
\usepackage{caption}
\examsetup{
  page/size=a4paper,
  paren/show-paren=true,
  paren/show-answer=true,
  fillin/show-answer=false,
  font      = times,
  math-font = xits,
  fillin/no-answer-type=none,
  %solution/show-solution=show-stay,
  %solution/label-indentation=false
  }
%\newcommand{\customcurve}[1][1]{%
  %\tikz[baseline,scale=#1,line width=0.8pt,line cap=round]{%
    %\draw (0.15,0.1) ellipse (0.15 and 0.1);
    %\draw (0,0.1) -- (-0.12 , -0.1);
    %\draw (0,0.1) -- (-0.12 , 0.3);
  %}%
%}
\newcommand{\customcurve}[1][1]{%
  \tikz[baseline,yshift=3,scale=0.05,thick,line width=0.8pt,
        line cap=round, domain=-1:2.832, samples=300]{
    \draw plot (\x,{sqrt(16/(\x+2)^2 - (\x-2)^2)});
    \draw plot (\x,{-sqrt(16/(\x+2)^2 - (\x-2)^2)});
    }
}

\everymath{\displaystyle}

\title{2026年04月05日 试卷类模板组卷2}

\subject{数学}



\begin{document}

% \information{
%   姓名\underline{\hspace{6em}},
%   座位号\underline{\hspace{15em}}
% }
% \warning{（在此卷上答题无效）}

\secret

\maketitle
\begin{center}
    本试卷共 \pageref{LastPage} 页，19 题。全卷满分 150 分。考试用时 120 分钟。
\end{center}


\begin{notice}
  \item 答题前，先将自己的姓名、准考证号、考场号、座位号填写在试卷和答题卡上，
    并将准考证号条形码粘贴在答题卡上的指定位置。
  \item 选择题的作答：每小题选出答案后，用 2B 铅笔把答题卡上对应题目的答案标号涂黑。
    写在试卷、草稿纸和答题卡上的非答题区域均无效。
  \item 填空题和解答题的作答：用黑色签字笔直接答在答题卡上对应的答题区域内。
    写在试卷、草稿纸和答题卡上的非答题区域均无效。
  \item 考试结束后，请将本试卷和答题卡一并上交。
\end{notice}



\section{
  选择题：本题共 8 小题，每小题 5 分，共 40 分。
  在每小题给出的四个选项中，只有一项是符合题目要求的。
}

% 1.
\begin{question}
已知 $z=-1-\mathrm{i}$，则 $|z|=$
(\hspace{1cm})
\begin{choices}
\item 0
\item 1
\item $\sqrt{2}$
\item 2
\end{choices}
\end{question}

% 2.
\begin{question}
设命题 $p: \forall x\in\mathbf{R}, |x+1|>1$；命题 $q: \exists x>0, x^3=x$. 则
(\hspace{1cm})
\begin{choices}
\item $p$ 和 $q$ 都是真命题
\item $\neg p$ 和 $q$ 都是真命题
\item $p$ 和 $\neg q$ 都是真命题
\item $\neg p$ 和 $\neg q$ 都是真命题
\end{choices}
\end{question}

% 3.
\begin{question}
已知向量 $\mathbf{a}, \mathbf{b}$ 满足 $|\mathbf{a}|=1$，$|\mathbf{a}+2\mathbf{b}|=2$，且 $(\mathbf{b}-2\mathbf{a})\perp\mathbf{b}$，则 $|\mathbf{b}|=$
(\hspace{1cm})
\begin{choices}
\item $\dfrac{1}{2}$
\item $\dfrac{\sqrt{2}}{2}$
\item $\dfrac{\sqrt{3}}{2}$
\item 1
\end{choices}
\end{question}

% 4.
\begin{question}
某农业研究部门在面积相等的 100 块稻田上种植一种新型水稻，得到各块稻田的亩产量（单位：kg）并整理如下表：
\begin{center}
\begin{tabular}{|l|c|c|c|c|c|c|}
\hline
亩产量 & [900,950) & [950,1000] & [1000,1050) & [1050,1100) & [1100,1150) & [1150,1200] \\
\hline
田块数 & 6 & 12 & 18 & 30 & 24 & 10 \\
\hline
\end{tabular}
\end{center}
据表中数据，下列结论正确的是
(\hspace{1cm})
\begin{choices}
\item 100 块稻田的亩产量的中位数小于 1050 kg
\item 100 块稻田的亩产量低于 1100 kg 的稻田所占比例超过 70\%
\item 100 块稻田的亩产量的极差介于 200 kg 到 300 kg 之间
\item 100 块稻田的亩产量的平均值介于 900 kg 到 1000 kg 之间
\end{choices}
\end{question}

% 5.
\begin{question}
已知曲线 $C: x^2+y^2=16(y>0)$，从 $C$ 上任意一点 $P$ 向 $x$ 轴作垂线 $PP'$，$P'$ 是垂足，则线段 $PP'$ 的中点的轨迹方程为
(\hspace{1cm})
\begin{choices}
\item $\dfrac{x^2}{16}+\dfrac{y^2}{4}=1(y>0)$
\item $\dfrac{x^2}{16}+\dfrac{y^2}{8}=1(y>0)$
\item $\dfrac{y^2}{16}+\dfrac{x^2}{4}=1(y>0)$
\item $\dfrac{y^2}{16}+\dfrac{x^2}{8}=1(y>0)$
\end{choices}
\end{question}

% 6.
\begin{question}
已知函数 $f(x)=\begin{cases}-x^2-2ax-a, & x<0 \\ e^x+\ln(x+1), & x\geqslant 0\end{cases}$ 在 $\mathbf{R}$ 上单调递增，则 $a$ 的取值范围是
(\hspace{1cm})
\begin{choices}
\item $(-\infty,0]$
\item $[-1,0]$
\item $[-1,1]$
\item $[0,+\infty)$
\end{choices}
\end{question}

% 7.
\begin{question}
已知正三棱台 $ABC-A_1B_1C_1$ 的体积为 $\dfrac{52}{3}$，$AB=6, A_1B_1=2$，则 $AA_1$ 与平面 $ABC$ 所成角的正切值为
(\hspace{1cm})
\begin{choices}
\item $\dfrac{1}{2}$
\item 1
\item 2
\item 3
\end{choices}
\end{question}

% 8.
\begin{question}
已知定义域为 $\mathbf{R}$ 的函数 $f(x)$ 满足 $f(x)>f(x-1)+f(x-2)$，且当 $x<3$ 时，$f(x)=x$，则
(\hspace{1cm})
\begin{choices}
\item $f(10)>100$
\item $f(20)>1000$
\item $f(10)<1000$
\item $f(20)<10000$
\end{choices}
\end{question}

% 9.
\begin{question}
对于函数 $f(x)=\sin 2x$ 和 $g(x)=\sin\left(2x-\dfrac{\pi}{4}\right)$，下列结论正确的有
(\hspace{1cm})
\begin{choices}
\item $f(x)$ 与 $g(x)$ 有相同零点
\item $f(x)$ 与 $g(x)$ 有相同最大值
\item $f(x)$ 与 $g(x)$ 有相同的最小正周期
\item $f(x)$ 与 $g(x)$ 的图像有相同的对称轴
\end{choices}
\end{question}

\section{
  选择题：本题共 3 小题，每小题 6 分，共 18 分。
  在每小题给出的选项中，有多项符合题目要求的。
  全部选对的得 6 分，部分选择的得部分分，有选错的得 0 分。
}

% 10.
\begin{question}
抛物线 $C: y^2=4x$ 的准线为 $l$，$P$ 为 $C$ 上的动点，过 $P$ 作 $\odot A: x^2+(y-4)^2=1$ 的一条切线，$Q$ 为切点. 过 $P$ 作 $l$ 的垂线，垂足为 $B$，则
(\hspace{1cm})
\begin{choices}
\item $l$ 与 $\odot A$ 相切
\item 当 $P,A,B$ 共线时，$|PQ|=\sqrt{15}$
\item 当 $|PB|=2$ 时，$PA\perp AB$
\item 满足 $|PA|=|PB|$ 的点 $P$ 有且仅有 2 个
\end{choices}
\end{question}

% 11.
\begin{question}
设函数 $f(x)=2x^3-3ax^2+1$，则
(\hspace{1cm})
\begin{choices}
\item 当 $a>1$ 时，$f(x)$ 有 3 个零点
\item 当 $a<0$ 时，$x=0$ 是 $f(x)$ 的极大值点
\item 存在 $a, b$，使得 $x=b$ 为曲线 $y=f(x)$ 的对称轴
\item 存在 $a$，使得点 $(1, f(1))$ 为曲线 $y=f(x)$ 的对称中心
\end{choices}
\end{question}

% 12.
\begin{question}
记 $S_n$ 为等差数列 $\{a_n\}$ 的前 $n$ 项和. 若 $a_3+a_4=7, 3a_2+a_5=5$，则 $S_{10}=$\underline{\hspace{4em}}.
\end{question}

\section{填空题：本题共 3 小题，每小题 5 分，共 15 分。}

% 13.
\begin{question}
已知 $\alpha$ 为第一象限的角，$\beta$ 为第三象限的角. 若 $\tan\alpha+\tan\beta=4, \tan\alpha\tan\beta=\sqrt{2}+1$，则 $\sin(\alpha+\beta)=$\underline{\hspace{4em}}.
\end{question}

% 14.
\begin{question}
在如图所示的 $4\times 4$ 方格表中选 4 个方格，要求每行和每列恰有一个方格被选中，则共有\underline{\hspace{4em}}种选法. 在所有符合上述要求的选法中，选中方格中的 4 个数之和的最大值是\underline{\hspace{4em}}.
\begin{center}
\begin{tabular}{|c|c|c|c|}
\hline
11 & 21 & 31 & 40 \\
\hline
12 & 22 & 33 & 42 \\
\hline
13 & 22 & 33 & 43 \\
\hline
15 & 24 & 34 & 44 \\
\hline
\end{tabular}
\end{center}
\end{question}

% 15.
\begin{question}
记 $\triangle ABC$ 的内角 $A,B,C$ 的对边分别为 $a,b,c$. 已知 $\sin A+\sqrt{3}\cos A=2$.

(1) 求 $A$;

(2) 若 $a=2, \sqrt{2}b\sin C=c\sin 2B$，求 $\triangle ABC$ 的周长.
\end{question}

\section{解答题：本题共 5 小题，共 77 分。解答应写出文字说明、证明过程或者演算步骤。}

% 16.
\begin{problem}[points = 15]
已知函数 $f(x)=e^x-ax-a^3$.

(1) 当 $a=1$ 时，求曲线 $y=f(x)$ 在点 $(1,f(1))$ 处的切线方程;

(2) 若 $f(x)$ 有极小值，且极小值小于 0，求 $a$ 的取值范围.
\end{problem}

% 17.
\begin{problem}[points = 15]
如图，平面四边形 $ABCD$ 中，$AB=8, CD=3, AD=5\sqrt{3}$，$\angle ADC=90^\circ, \angle BAD=30^\circ$，点 $E,F$ 满足 $\overrightarrow{AE}=\dfrac{2}{5}\overrightarrow{AD}, \overrightarrow{AF}=\dfrac{1}{2}\overrightarrow{AB}$. 将 $\triangle AEF$ 沿 $EF$ 对折至 $\triangle PEF$，使得 $PC=4\sqrt{3}$.

(1) 证明：$EF\perp PD$;

(2) 求面 $PCD$ 与面 $PBF$ 所成二面角的正弦值.
\end{problem}

% 18.
\begin{problem}[points = 17]
某投篮比赛分为两个阶段，每个参赛队由两名队员组成。比赛规则如下：第一阶段由参赛队中一名队员投篮三次。若三次都未投中，则该队被淘汰，比赛成绩为 $0$ 分；若至少投中一次，则该队进入第二阶段，由该队的另一名队员投篮 $3$ 次，每次投中得 $5$ 分，未投中得 $0$ 分。该队的比赛成绩为第二阶段的得分总和。某参赛队由甲、乙两名队员组成，设甲每次投中的概率为 $p$，乙每次投中的概率为 $q$，各次投中与否相互独立。

（1）若 $p=0.4$，$q=0.5$，甲参加第一阶段比赛，求甲、乙所在队的比赛成绩不少于 $5$ 分的概率；

（2）假设 $0<p<q$.

（i）为使得甲、乙所在队的比赛成绩为 $15$ 分的概率最大，应该由谁参加第一阶段比赛？

（ii）为使得甲、乙所在队的比赛成绩的数学期望最大，应该由谁参加第一阶段比赛？
\end{problem}

% 19.
\begin{problem}[points = 17]
已知双曲线 $C: x^2-y^2=m(m>0)$，点 $P_1(5,4)$ 在 $C$ 上. $k$ 为常数，$0<k<1$. 按照如下方法依次构造点 $P_n(n=2,3,\cdots)$：过 $P_{n-1}$ 作斜率为 $k$ 的直线与 $C$ 的左支交于点 $Q_{n-1}$，令 $P_n$ 为 $Q_{n-1}$ 关于 $y$ 轴的对称点. 记 $P_n$ 的坐标为 $(x_n, y_n)$.

（1）若 $k=\dfrac{1}{2}$，求 $x_2, y_2$;

（2）证明：数列 $\{x_n-y_n\}$ 是公比为 $\dfrac{1+k}{1-k}$ 的等比数列;

（3）记 $S_n$ 为 $\triangle P_n P_{n+1} P_{n+2}$ 的面积，证明：对任意正整数 $n$，$S_n=S_{n+1}$.
\end{problem}


\end{document}